\chapter{Problem Formulation}
\label{chap:formulation}

\vspace*{\fill}

This chapter introduces the class of problems considered throughout this work: linear
programs with a multi-block structure, whose blocks are coupled by linking constraints that
may be equalities or inequalities.

\medskip

We first state the problem in its original (compact) form and then apply the
\textbf{Dantzig--Wolfe reformulation}, which replaces the variables of each block by convex
combinations of the extreme points of the block. This leads to the \textbf{master problem},
a linear program with a very large number of columns, which is solved by
\textbf{column generation}. We define the restricted master problem, the reduced cost, the
pricing problem and the Lagrangian bound, and collect them in the classical column
generation algorithm.

\medskip

Finally, we introduce \textbf{surplus variables} for the inequality rows. They bring every
linking constraint into equality form and lead to a single, \emph{unified} formulation of the
master problem in which equality and inequality rows are treated alike. This formulation is
the starting point for everything that follows. It also reveals the one structural property
that the inequality rows destroy: the surplus variables belong to no block.

\vspace*{\fill}

\newpage


\section{Multi-Block Linear Programs with Mixed Linking Constraints}
\label{sec:multiblock}

We consider a linear program whose variables are partitioned into $K$ blocks
$x_1, \dots, x_K$, with $x_k \in \R^{n_k}$. Each block has its own constraints, and the blocks
interact only through a set of linking constraints.

\begin{definition}[Multi-block linear program] \label{def:multiblock}
A \textbf{multi-block linear program with mixed linking constraints} is a problem of the
form
\begin{subequations} \label{eq:compact}
\begin{align}
\min_{x_1,\dots,x_K} \quad & \sum_{k=1}^{K} c_k^\top x_k \\
\text{s.t.} \quad
& \sum_{k=1}^{K} A_{E,k}\, x_k = b_E, \label{eq:compact-eq} \\
& \sum_{k=1}^{K} A_{I,k}\, x_k \ge b_I, \label{eq:compact-ineq} \\
& x_k \in P_k, \qquad k = 1,\dots,K, \label{eq:compact-block}
\end{align}
\end{subequations}
where $A_{E,k} \in \R^{m_E\times n_k}$ and $A_{I,k}\in\R^{m_I\times n_k}$ are the
coefficients of block $k$ in the $m_E$ \textbf{equality linking rows} and the $m_I$
\textbf{inequality linking rows}, and
\[
P_k = \{\, x_k \in \R^{n_k} \;:\; F_k x_k = d_k,\; l_k \le x_k \le u_k \,\}
\]
is the \textbf{block polytope} of block $k$. We write $m = m_E + m_I$ for the total number of
linking rows. Constraints~\eqref{eq:compact-eq}--\eqref{eq:compact-ineq} are the
\textbf{linking constraints}, and constraints~\eqref{eq:compact-block} are the
\textbf{block constraints}.
\end{definition}

\begin{assumption} \label{ass:bounded}
Every block polytope $P_k$ is nonempty and bounded, and problem~\eqref{eq:compact} is
feasible.
\end{assumption}

Under \cref{ass:bounded} the feasible region of~\eqref{eq:compact} is a nonempty polytope,
so the problem has an optimal solution. Its optimal value is denoted by $z^\star$.

\bigskip

The constraint matrix of~\eqref{eq:compact} has the \textbf{block-angular} structure shown
in \cref{fig:block-angular}. The block constraints form a block-diagonal matrix, and only
the linking rows couple different blocks. If the linking rows were absent, the problem would
decompose into $K$ independent problems over $P_1, \dots, P_K$. Column generation exploits
exactly this observation.

\begin{figure}[ht]
\centering
\begin{tikzpicture}[x=0.55cm, y=0.55cm]
    % linking rows
    \fill[cblue!25] (0,10) rectangle (16,12);
    \fill[corange!35] (0,8.6) rectangle (16,10);
    \foreach \k in {0,1,2,3} {
        \draw[cblue!70!black] (4*\k,10) rectangle (4*\k+4,12);
        \draw[corange!80!black] (4*\k,8.6) rectangle (4*\k+4,10);
        \node at (4*\k+2,11) {\small $A_{E,\the\numexpr\k+1\relax}$};
        \node at (4*\k+2,9.3) {\small $A_{I,\the\numexpr\k+1\relax}$};
    }
    % block-diagonal part
    \foreach \k in {0,1,2,3} {
        \fill[cgreen!25] (4*\k, 8.6-2*\k-2) rectangle (4*\k+4, 8.6-2*\k);
        \draw[cgreen!70!black] (4*\k, 8.6-2*\k-2) rectangle (4*\k+4, 8.6-2*\k);
        \node at (4*\k+2, 8.6-2*\k-1) {\small $F_{\the\numexpr\k+1\relax}$};
    }
    \draw[thick] (0,0.6) rectangle (16,12);
    % labels
    \draw[decorate, decoration={brace, amplitude=4pt}] (16.3,12) -- (16.3,10)
        node[midway, right=5pt] {\small equality rows, $=b_E$};
    \draw[decorate, decoration={brace, amplitude=4pt}] (16.3,10) -- (16.3,8.6)
        node[midway, right=5pt] {\small inequality rows, $\ge b_I$};
    \draw[decorate, decoration={brace, amplitude=4pt}] (16.3,8.6) -- (16.3,0.6)
        node[midway, right=5pt] {\small block constraints, $x_k\in P_k$};
    \foreach \k in {0,1,2,3}
        \node at (4*\k+2, 12.6) {\small $x_{\the\numexpr\k+1\relax}$};
\end{tikzpicture}
\caption{Block-angular structure of the multi-block linear program~\eqref{eq:compact} for
$K = 4$ blocks. The linking rows (top) couple all blocks; the block constraints (bottom) are
block-diagonal.}
\label{fig:block-angular}
\end{figure}

\begin{remark}[Inequalities of the other direction] \label{rem:le-rows}
A linking row of the form $\sum_k a_{i,k}^\top x_k \le b_i$ is the row
$\sum_k (-a_{i,k})^\top x_k \ge -b_i$, so the form~\eqref{eq:compact-ineq} covers both
directions. For the presentation we therefore only consider ``$\ge$'' rows. In
\cref{sec:unified} we show that a ``$\le$'' row can equivalently be kept as it is, with a
surplus variable of the opposite sign. This is how such rows are handled in the
implementation.
\end{remark}

\begin{example}[Crew scheduling] \label{ex:crew}
In the crew scheduling problem studied in \cref{chap:results}, block $k$ is the schedule of
pilot $k$: a path from a source to a sink in a time--space network, with $F_k$ the
node--edge incidence matrix of the network. The linking rows are the \emph{coverage rows}:
row $i$ counts how often flight $i$ is operated, summed over all pilots. An equality row
``$=1$'' requires that the flight is operated exactly once, and an inequality row ``$\ge 1$''
requires it at least once, which allows a pilot to travel on a flight as a passenger.
\end{example}


\section{Dantzig--Wolfe Reformulation}
\label{sec:dantzig-wolfe}

Since $P_k$ is a nonempty polytope, it is the convex hull of its finitely many extreme points
(Minkowski's theorem). Let $\mathcal{J}_k$ be an index set of the extreme points of $P_k$,
and let $x^{(j)}$, $j\in\mathcal{J}_k$, denote these extreme points. Every $x_k\in P_k$ can be
written as
\begin{equation} \label{eq:convex-combination}
x_k = \sum_{j\in\mathcal{J}_k} \lambda_j\, x^{(j)}, \qquad
\sum_{j\in\mathcal{J}_k}\lambda_j = 1, \qquad \lambda_j \ge 0 .
\end{equation}
Substituting~\eqref{eq:convex-combination} into~\eqref{eq:compact} replaces the variables
$x_k$ by the weights $\lambda_j$. Each extreme point turns into one column of a new problem,
which motivates the following definitions.

\begin{definition}[Columns] \label{def:columns}
Let $\mathcal{J} = \mathcal{J}_1\cup\dots\cup\mathcal{J}_K$ (disjoint union). For
$j\in\mathcal{J}_k$ we define
\begin{itemize}
    \item the \textbf{block} of column $j$: $\psi(j) = k$;
    \item the \textbf{cost} of column $j$: $c_j = c_k^\top x^{(j)}$;
    \item the \textbf{linking coefficients} of column $j$:
          $a_j^E = A_{E,k}\,x^{(j)}\in\R^{m_E}$, $a_j^I = A_{I,k}\,x^{(j)}\in\R^{m_I}$, and
          $a_j = \bigl[\,a_j^E;\, a_j^I\,\bigr]\in\R^{m}$.
\end{itemize}
We collect them in the cost vector $c = (c_j)_{j\in\mathcal{J}}$ and the matrices
\[
A_E = \bigl[\, a_j^E \,\bigr]_{j\in\mathcal{J}}\in\R^{m_E\times|\mathcal{J}|}, \qquad
A_I = \bigl[\, a_j^I \,\bigr]_{j\in\mathcal{J}}\in\R^{m_I\times|\mathcal{J}|}, \qquad
A = \begin{bmatrix} A_E \\ A_I\end{bmatrix}\in\R^{m\times|\mathcal{J}|}.
\]
\end{definition}

\begin{definition}[Convexity matrix] \label{def:convexity}
The \textbf{convexity matrix} $C\in\{0,1\}^{K\times|\mathcal{J}|}$ has the entries
$C_{kj} = 1$ if $\psi(j) = k$ and $C_{kj} = 0$ otherwise. The constraint $C\lambda = \one$
states that the weights of the columns of every block sum to one.
\end{definition}

\begin{definition}[Master problem] \label{def:master}
The \textbf{master problem} is the linear program
\begin{subequations} \label{eq:master}
\begin{align}
\min_{\lambda} \quad & c^\top \lambda \\
\text{s.t.} \quad & A_E\,\lambda = b_E, && (\pi_E) \label{eq:master-eq}\\
& A_I\,\lambda \ge b_I, && (\pi_I) \label{eq:master-ineq}\\
& C\lambda = \one, && (\mu) \label{eq:master-conv}\\
& \lambda \ge 0, \label{eq:master-nonneg}
\end{align}
\end{subequations}
where the symbols in parentheses denote the dual variables of the constraints.
\end{definition}

\begin{theorem}[Dantzig--Wolfe \cite{dantzig1960decomposition}] \label{thm:dw}
Under \cref{ass:bounded}, problem~\eqref{eq:master} has the optimal value $z^\star$. If
$\lambda$ is optimal for~\eqref{eq:master}, then $x_k = \sum_{j\in\mathcal{J}_k}\lambda_j\,
x^{(j)}$, $k=1,\dots,K$, is optimal for~\eqref{eq:compact}.
\end{theorem}

\begin{proof}
By~\eqref{eq:convex-combination}, $x_k\in P_k$ if and only if $x_k$ is a convex combination
of the $x^{(j)}$, $j\in\mathcal{J}_k$. Under this substitution the objective and the
linking rows are linear in $\lambda$, with exactly the coefficients of \cref{def:columns}.
The feasible points of~\eqref{eq:compact} are therefore the images of the feasible points
of~\eqref{eq:master}, with equal objective values. 
\end{proof}

The convexity matrix has a very particular structure, which will be essential for the
nullspace method of \cref{chap:nullspace}.

\begin{definition}[Column generation property] \label{def:cg-property}
A matrix $C\in\R^{K\times n}$ has the \textbf{column generation property} if
\begin{enumerate}
    \item every entry of $C$ is $0$ or $1$;
    \item every column of $C$ contains exactly one entry equal to $1$;
    \item every row of $C$ contains at least one entry equal to $1$.
\end{enumerate}
\end{definition}

\begin{lemma} \label{lem:cg-property}
The convexity matrix $C$ of \cref{def:convexity} has the column generation property, and so
does every submatrix of $C$ formed by a subset of its columns that contains at least one
column of every block.
\end{lemma}

\begin{proof}
Property~(i) holds by definition. Every column belongs to exactly one block, which gives
property~(ii) for every subset of columns. Property~(iii) holds for a subset of columns if
and only if every block is represented in it. 
\end{proof}

\begin{lemma} \label{lem:C-rank}
If $C\in\R^{K\times n}$ has the column generation property, then $\rank(C) = K$.
\end{lemma}

\begin{proof}
By property~(iii), for every row $k$ there is a column $j_k$ with $C_{kj_k} = 1$. By
property~(ii), this column has no other nonzero entry, i.e.\ it equals the unit vector
$e_k$. The columns $j_1,\dots,j_K$ therefore form the identity matrix $I_K$. 
\end{proof}

\subsection*{The dual of the master problem}

The dual of~\eqref{eq:master} has one constraint per column and the sign condition of the
inequality rows:
\begin{subequations} \label{eq:master-dual}
\begin{align}
\max_{\pi_E,\,\pi_I,\,\mu} \quad & b_E^\top \pi_E + b_I^\top\pi_I + \one^\top\mu \\
\text{s.t.} \quad & c_j - a_j^\top \pi - \mu_{\psi(j)} \ge 0, \qquad j\in\mathcal{J},
\label{eq:dual-columns}\\
& \pi_I \ge 0, \label{eq:dual-sign}
\end{align}
\end{subequations}
where $\pi = [\pi_E;\pi_I]\in\R^m$. Compared with a problem that has only equality linking
rows, the inequality rows add the \textbf{sign condition}~\eqref{eq:dual-sign}. This is the
only change in the dual, but much of the work in the following chapters is devoted to it.


\section{Column Generation}
\label{sec:column-generation}

The index set $\mathcal{J}$ is typically far too large to enumerate. In crew scheduling, for
example, $\mathcal{J}_k$ contains every feasible schedule of pilot $k$. Column generation
solves~\eqref{eq:master} by working with a small subset of the columns at a time.

\begin{definition}[Restricted master problem] \label{def:rmp}
For a subset $\mathcal{J}'\subseteq\mathcal{J}$ that contains at least one column of every
block, the \textbf{restricted master problem} is problem~\eqref{eq:master} with all columns
outside $\mathcal{J}'$ removed (equivalently, with $\lambda_j = 0$ fixed for
$j\notin\mathcal{J}'$). We write $n = |\mathcal{J}'|$ for the number of columns in the
restricted master problem.
\end{definition}

\begin{definition}[Reduced cost] \label{def:reduced-cost}
For dual values $(\pi,\mu)$ the \textbf{reduced cost} of a column $j\in\mathcal{J}$ is
\[
\bar c_j(\pi,\mu) \;=\; c_j - a_j^\top\pi - \mu_{\psi(j)} .
\]
\end{definition}

The reduced cost is the slack of the dual constraint~\eqref{eq:dual-columns}. A column with
a negative reduced cost violates the dual constraints and can improve the restricted master
problem. Finding the column of most negative reduced cost in block $k$ is the
\textbf{pricing problem}
\begin{equation} \label{eq:pricing}
\min_{j\in\mathcal{J}_k}\bar c_j(\pi,\mu)
\;=\; \min_{x_k\in P_k}\bigl(c_k - A_k^\top\pi\bigr)^\top x_k \;-\;\mu_k,
\qquad A_k = \begin{bmatrix}A_{E,k}\\ A_{I,k}\end{bmatrix},
\end{equation}
where the equality holds because a linear function attains its minimum over a polytope at an
extreme point. The pricing problem is a problem over a single block. In crew scheduling it
is a shortest path problem with edge weights $c_{e,k} - (A_k^\top\pi)_e$.

\begin{prop}[Optimality test] \label{prop:cg-optimality}
Let $\lambda'$ be optimal for the restricted master problem, with optimal duals
$(\pi,\mu)$. If $\pi_I\ge0$ and the pricing problem~\eqref{eq:pricing} has a nonnegative
optimal value for every block, then $\lambda'$, extended by zeros, is optimal for the master
problem~\eqref{eq:master}.
\end{prop}

\begin{proof}
The extended $\lambda'$ is feasible for~\eqref{eq:master}. The pair $(\pi,\mu)$ satisfies
the sign condition~\eqref{eq:dual-sign} by assumption and every dual
constraint~\eqref{eq:dual-columns} by the pricing result, so it is dual feasible. Complementary
slackness holds for the columns in $\mathcal{J}'$ by optimality of the restricted problem,
and trivially for the zero weights outside $\mathcal{J}'$. 
\end{proof}

For a linear restricted master problem the duals satisfy $\pi_I\ge0$ automatically, so the
pricing test alone decides optimality. With the regularized master problems of the next
chapter this will no longer be the case. Column generation also provides a lower bound on
$z^\star$ at every iteration in which all blocks are priced.

\begin{prop}[Lagrangian bound] \label{prop:lagrangian-bound}
For every $\pi = [\pi_E;\pi_I]$ with $\pi_I\ge0$, the \textbf{Lagrangian bound}
\[
\theta(\pi) \;=\; b^\top\pi + \sum_{k=1}^{K}\,\min_{x_k\in P_k}\bigl(c_k - A_k^\top\pi\bigr)^\top x_k,
\qquad b = \begin{bmatrix}b_E\\ b_I\end{bmatrix},
\]
satisfies $\theta(\pi)\le z^\star$.
\end{prop}

\begin{proof}
Let $x^\star$ be optimal for~\eqref{eq:compact}. Because $\sum_k A_{E,k}x^\star_k = b_E$,
$\sum_k A_{I,k}x^\star_k\ge b_I$ and $\pi_I\ge0$, we have
$\pi^\top\bigl(b - \sum_k A_k x_k^\star\bigr)\le 0$. Hence
\[
z^\star = \sum_k c_k^\top x^\star_k
\;\ge\; \sum_k c_k^\top x^\star_k + \pi^\top\Bigl(b - \sum_k A_k x_k^\star\Bigr)
= b^\top\pi + \sum_k\bigl(c_k - A_k^\top\pi\bigr)^\top x^\star_k
\;\ge\; \theta(\pi),
\]
since $x_k^\star\in P_k$. 
\end{proof}

\begin{remark} \label{rem:bound-sign}
The sign condition $\pi_I\ge0$ in \cref{prop:lagrangian-bound} is essential: for $\pi_i<0$
in an inequality row, the term $\pi_i\,(b_i - \sum_k a_{i,k}^\top x_k)$ can become arbitrarily
large. In \cref{sec:center-update} the bound is therefore evaluated at the projection of
the dual vector onto $\{\pi_I\ge0\}$.
\end{remark}

The resulting method is summarized in \cref{alg:classical-cg}. The initial columns are
usually obtained by solving the pricing problem once with zero duals, which gives one
column per block.

\begin{algobox}{Classical Column Generation}{alg:classical-cg}
\State \textbf{Input:} initial columns $\mathcal{J}'$ with at least one column per block
\Repeat
    \State Solve the (linear) restricted master problem over $\mathcal{J}'$; obtain
           $\lambda'$ and duals $(\pi,\mu)$
    \For{each block $k = 1,\dots,K$}
        \State Solve the pricing problem~\eqref{eq:pricing}; let $j$ be the best column
        \If{$\bar c_j(\pi,\mu) < 0$}
            \State $\mathcal{J}'\gets\mathcal{J}'\cup\{j\}$
        \EndIf
    \EndFor
\Until{no column with negative reduced cost was found}
\State \textbf{Output:} $\lambda'$, optimal for the master problem by \cref{prop:cg-optimality}
\end{algobox}

\begin{remark}[Instability] \label{rem:instability}
The duals of a linear restricted master problem are extreme points of its dual feasible
region. When a column is added, the dual solution can jump to a very different extreme point,
even if the primal solution hardly changes. As a result, the duals oscillate strongly from
one iteration to the next, the pricing problem generates many columns that are never used,
and the method converges slowly in its final phase (the \emph{tailing-off effect}). In the
dual space, column generation is Kelley's cutting-plane method
\cite{kelley1960cutting} applied to the Lagrangian dual function, a method known for exactly
this behaviour. This is the motivation for the regularization of \cref{chap:regularization}.
\end{remark}


\section{Surplus Variables and the Unified Formulation}
\label{sec:unified}

The two kinds of linking rows in~\eqref{eq:master} behave differently: the equality rows
fix $A_E\lambda$, while the inequality rows only bound $A_I\lambda$ from below. To treat both
in the same way, we introduce for every inequality row $i$ a \textbf{surplus variable}
$s_i\ge0$ with
\[
(A_I\lambda)_i - s_i = (b_I)_i, \qquad\text{i.e.}\qquad A_I\lambda - s = b_I .
\]
The surplus measures how much the row is over-satisfied.

\begin{definition}[Unified formulation] \label{def:unified}
Let $s\in\R^{m_I}$ be the vector of surplus variables. The \textbf{unified variable} and the
\textbf{augmented data} are
\[
y = \begin{bmatrix}\lambda\\ s\end{bmatrix}, \qquad
\Abar = \begin{bmatrix} A_E & 0\\ A_I & -I\end{bmatrix}, \qquad
\Cbar = \begin{bmatrix} C & 0\end{bmatrix}, \qquad
\cbar = \begin{bmatrix} c\\ 0\end{bmatrix}, \qquad
b = \begin{bmatrix}b_E\\ b_I\end{bmatrix}.
\]
The \textbf{unified master problem} is
\begin{equation} \label{eq:unified}
\min_{y}\;\cbar^\top y \qquad\text{s.t.}\qquad \Abar y = b,\quad \Cbar y = \one,\quad y\ge0 .
\end{equation}
\end{definition}

\begin{lemma} \label{lem:unified}
Problems~\eqref{eq:master} and~\eqref{eq:unified} are equivalent: $\lambda$ is feasible
for~\eqref{eq:master} if and only if $y = [\lambda;\, A_I\lambda - b_I]$ is feasible
for~\eqref{eq:unified}, with the same objective value. Moreover, the dual
of~\eqref{eq:unified} is the dual~\eqref{eq:master-dual} of the master problem.
\end{lemma}

\begin{proof}
The first statement follows from $s = A_I\lambda - b_I\ge0$. For the second, the dual
of~\eqref{eq:unified} has multipliers $\pi$ for $\Abar y = b$ and $\mu$ for $\Cbar y = \one$,
and one constraint per column of $[\Abar;\Cbar]$. For the columns of $\lambda$ these are the
constraints~\eqref{eq:dual-columns}. The column of surplus $s_i$ has cost $0$, coefficient
$-e_i$ in the inequality rows and $0$ in the convexity rows, so its dual constraint reads
$0 - (-e_i)^\top\pi_I - 0 \ge 0$, i.e.\ $(\pi_I)_i\ge0$. These are exactly the sign
conditions~\eqref{eq:dual-sign}. 
\end{proof}

\Cref{lem:unified} shows that the surplus columns \emph{encode the sign conditions as
ordinary dual constraints}. Problem~\eqref{eq:unified} has the same form as a master problem
with equality linking rows only, with $(A, C, c)$ replaced by $(\Abar,\Cbar,\cbar)$. All the
following chapters work with this form.

\begin{remark}[The column generation property is lost] \label{rem:cg-property-lost}
The augmented convexity matrix $\Cbar = [\,C\;\;0\,]$ has $m_I$ zero columns, one per
surplus variable, so it violates property~(ii) of \cref{def:cg-property}. The surpluses do
not belong to any block. \Cref{chap:qpas} shows that this causes no difficulty: zero
columns are even simpler to handle than unit columns, but they have to be treated as a
separate group of variables.
\end{remark}

\begin{remark}[``$\le$'' rows and the slack sign] \label{rem:kappa}
A row $a_i^\top\lambda\le b_i$ can be written with a \emph{slack} $s_i\ge0$ as
$a_i^\top\lambda + s_i = b_i$. In general, let $\kappa_i\in\{-1,+1\}$ be the coefficient of
the slack of inequality row $i$ ($\kappa_i = -1$ for a ``$\ge$'' row, $\kappa_i = +1$ for a
``$\le$'' row) and $K_s = \diag(\kappa)$. The unified formulation then uses
\[
\Abar = \begin{bmatrix} A_E & 0\\ A_I & K_s\end{bmatrix}
\]
instead of \cref{def:unified}, and the dual constraint of slack $i$ becomes
$-\kappa_i(\pi_I)_i\ge0$, i.e.\ $(\pi_I)_i\le0$ for a ``$\le$'' row. Since
$\kappa_i^2 = 1$, every result of the following chapters holds for both kinds of rows; only a
few formulas carry a factor $\kappa_i$, and these are collected in \cref{rem:kappa-summary}.
In the text we write all formulas for ``$\ge$'' rows ($K_s = -I$).
\end{remark}

\begin{remark}[Range rows] \label{rem:range}
A row with two different finite bounds $\ell_i\le a_i^\top\lambda\le u_i$ needs a slack with
an upper bound. Such rows are not considered in this work.
\end{remark}


\section{Summary of Notation}

\begin{table}[ht]
\centering
\renewcommand{\arraystretch}{1.2}
\begin{tabular}{ll}
\toprule
symbol & meaning \\
\midrule
$K$, $k$ & number of blocks, block index \\
$m_E$, $m_I$, $m$ & number of equality, inequality and all linking rows \\
$\mathcal{J}$, $\mathcal{J}_k$, $\mathcal{J}'$ & all columns, columns of block $k$, columns of the restricted problem \\
$n$ & number of columns in the restricted master problem \\
$\psi(j)$ & block of column $j$ \\
$c_j$, $a_j = [a_j^E; a_j^I]$ & cost and linking coefficients of column $j$ \\
$A_E$, $A_I$, $A$ & linking coefficient matrices of the columns \\
$C$ & convexity matrix \\
$\lambda$, $s$, $y = [\lambda; s]$ & column weights, surpluses, unified variable \\
$\Abar$, $\Cbar$, $\cbar$ & augmented data of the unified formulation \\
$\pi = [\pi_E;\pi_I]$, $\mu$ & duals of the linking rows and of the convexity rows \\
$\bar c_j(\pi,\mu)$ & reduced cost of column $j$ \\
$\theta(\pi)$ & Lagrangian bound \\
\bottomrule
\end{tabular}
\caption{Notation of \cref{chap:formulation}.}
\end{table}
